\documentclass[11pt,a4paper]{article}
% Shared preamble for RuPhO English problem statements (match T1 2026 style).
% Use: \input{latex/rupho_en_preamble.tex} after \documentclass.
\usepackage[margin=1in]{geometry}
\usepackage{amsmath,amssymb}
\usepackage{graphicx}
\usepackage{float}
\usepackage{enumitem}
\usepackage{microtype}
\usepackage{caption}
\usepackage{tabularx}
\usepackage{hyperref}

\captionsetup{font=small,labelfont=bf,skip=6pt}

\setlist[itemize]{leftmargin=1.2cm}
\setlist[enumerate]{leftmargin=1.2cm}

% One outer box per subproblem: [id | statement | points] (no inner rules)
\newcommand{\problembox}[3]{%
  \par\addvspace{0.42em}%
  \noindent
  \setlength{\fboxsep}{7.5pt}%
  \setlength{\fboxrule}{0.95pt}%
  \fbox{%
    \begin{minipage}{\dimexpr\linewidth-2\fboxsep-2\fboxrule\relax}
    \begin{tabularx}{\linewidth}{@{}>{\raggedright\arraybackslash\bfseries}p{2.1em} @{\hspace{0.9em}} >{\raggedright\arraybackslash}X @{\hspace{0.9em}} >{\raggedleft\arraybackslash\bfseries}p{2.4em}@{}}
    #1 & #2 & #3
    \end{tabularx}
    \end{minipage}%
  }%
  \par\addvspace{0.32em}%
}


\title{\textbf{Accelerated observers in the service station}}
\author{\normalsize Y25 (2025) --- English translation}
\date{}
\begin{document}
\maketitle

Typically, the special theory of relativity (SRT) considers motion only relative to inertial reference systems. However, the theory also considers bodies moving with acceleration, which can also be considered as observers making measurements. By connecting the reference system with accelerated observers, it is possible to obtain analogues of non-inertial reference systems in STR. Possible definitions of accelerated motion and the way in which an accelerated observer determines coordinates are not unambiguous. In this problem we will look at some of these constructions. For simplicity, in the entire problem we will consider only one-dimensional motion along the $Ox$ axis. Unless otherwise stated, the positions of observers are specified relative to some inertial frame of reference (CO) with coordinate $x$ and time $t$. As usual in SRT, we will use the concept of an event - any physical process occurring at a given point at a given moment in time. In a given reference frame in the one-dimensional case, an event is characterized by time $t$ and coordinate $x$, which for brevity we will sometimes call event coordinates.

\section*{Part A. Pseudo-relativistic accelerated observers (3.2 points)}

In this part we will not consider any relativistic effects, except for the finiteness of the signal propagation speed. Let's consider a system of initially stationary observers (point particles) that move with constant acceleration in the non-relativistic sense (that is, with a constant derivative of speed with respect to time in a given coordinate system). Observers can exchange signals that travel at the speed of light $c$. The problem is one-dimensional; we consider only the movement along the $x$ axis. All coordinates and times in this part are determined relative to a fixed inertial reference frame. The observer's acceleration depends on his initial coordinate $d$ as $$ a(d) = \textbackslash\{\}frac\{a_0\}\{1+ 2 a_0 d/c^2\}, \textbackslash\{\}quad d \textbackslash\{\}ge 0. $$ Such motion cannot continue indefinitely, since particle speeds cannot exceed the speed of light.

\problembox{A1}{Find the maximum possible movement time $t_{\max}$. In what follows, we will only consider motion at times not exceeding $t_{\max}$.}{0.10}

\problembox{A2}{Show that if for two observers at the initial moment $d_1 < d_2$, then at $t< t_{\max}$ their coordinates will still hold $x_1(t) < x_2(t)$.}{0.50}

\problembox{A3}{Let observer 1 at time $t_1$ emit a light signal, and observer 2 receive it at time $t_2$. Accelerations of observers $a_1$ and $a_2$, initial coordinates $d_1$ and $d_2$ respectively. Show that between these times the relation $$ \textbackslash\{\}left(t_1 - \textbackslash\{\}frac\{c\}\{a_1\} \textbackslash\{\}right)^2 = \textbackslash\{\}varepsilon \textbackslash\{\}left( t_2 - \textbackslash\{\}frac\{c\}\{a_2\}\textbackslash\{\}right)^2, $$если $d_2 > d_1$, и соотношение $$ \textbackslash\{\}left(t_1 + \textbackslash\{\}frac\{c\}\{a_1\} \textbackslash\{\}right)^2 =  \textbackslash\{\}varepsilon\textbackslash\{\}left( t_2 + \textbackslash\{\}frac\{c\}\{a_2\}\textbackslash\{\}right)^2, $$если $d_2 < d_1$. Выразите постоянную $\varepsilon$ через $a_1$ и $a_2$ holds.}{1.00}

\problembox{A4}{In order to measure the distance to another observer, observer 1 (with the initial coordinate value $d_1$) emits a ray at time $t_A$, which reaches observer 2 (with the initial coordinate $d_2> d_1$) at some time $t_C$, is reflected and returns to the first observer at time $t_B$ (according to the laboratory clock). Then observer 1 counts. that the distance to observer 2 is $D_{rad} = c (t_B- t_A)/2$. Express this distance in terms of $d_1$, $d_2$, $a_0$, $c$ and show that it is independent of $t_A$.}{1.30}

\problembox{A5}{Let $d_1 = d$ and $d_2 = d_1 + \Delta d$, where $\Delta d \ll d$ is the positive distance between observers. Express $D_{rad}$ to first order of smallness in $\Delta d$ in terms of $d$, $a_0$, $c$ and $\Delta d$.}{0.30}

Thus, we saw that due to the finite speed of propagation of signals, the distance to each other measured by observers turns out to be constant, although they move with different accelerations, and the distance in the laboratory CO changes.

\section*{Part B. Radar coordinates and Lorentz transformations (1.3 points)}

Let an observer move in some inertial reference frame with coordinates $(t, x)$, whose coordinate and time depend on his own time as $(t_O(\tau), x_O(\tau))$.  (Proper time is the time according to the watch of the observer himself.) In order to determine the coordinates and times of events, this observer illuminates everything around with radar beams (radar radiation moves at a constant speed $c$) and measures the times when the reflected signal reaches it. Let some event occur at time $t$, at a point with coordinate $x$. Let $\tau_1$ be the proper time at which the observer had to send a signal for this signal to reach the event, and $\tau_2$ be the proper time at which the observer would receive the signal emitted at the moment of the event. Then he assigns to the event the time and coordinate $$ T = \textbackslash\{\}frac\{1\}\{2\} (\textbackslash\{\}tau_1 + \textbackslash\{\}tau_2), \textbackslash\{\}qquad X = \textbackslash\{\}pm\textbackslash\{\}frac\{c\}\{2\} (\textbackslash\{\}tau_2 - \textbackslash\{\}tau_1).$$Знаку плюс отвечает случай, когда событие произошло справа от наблюдателя (при больших значениях координаты $x$), the minus sign is on the left.  In this part we will check that for an observer moving at a constant speed, this design reproduces standard coordinates in a moving inertial reference frame.

\problembox{B1}{Let the observer move relative to the inertial reference frame with speed $v$. Moreover, at the moment of time $t = 0$, $x = 0$. Express $x$, $t$ in terms of $\tau$.}{0.20}

\problembox{B2}{Let some event occur at a point with time and coordinate $(t, x)$. Determine the proper time values ​​$\tau_1$, $\tau_2$ for it. For simplicity, you can carry out calculations only for the case when the event occurred to the right of the observer.}{0.70}

\problembox{B3}{Show that the time and coordinate $(T, X)$ determined in this way coincide with the time and coordinate in the inertial FR moving with speed $v$ relative to the original FR.}{0.40}

\section*{Part C: Accelerated Observers (5.5 points)}

Now consider an observer moving with constant acceleration $g$ in an accompanying reference frame (an inertial reference frame in which the observer’s speed is zero at a given moment in time).

\problembox{C1}{Show that, with a suitable choice of initial conditions for such an observer, the dependence of the coordinate and time on the proper time has the form $$ x = A \textbackslash\{\}operatorname\{ch\} \textbackslash\{\}alpha\textbackslash\{\}tau, \textbackslash\{\}qquad t = B \textbackslash\{\}operatorname\{sh\}\textbackslash\{\}alpha \textbackslash\{\}tau. $$Определите значения постоянных $A$, $B$, $\alpha$. Выразите ответ через $g$, $c$.}{0.80}

\problembox{C2}{Let at the moment of time $t$ at a point with coordinate $x$ some event occur. These coordinates $(t, x)$ can be expressed in terms of $X$ and $T$ (defined for the observer from $\mathbf{C1}$ in the manner described in the introduction to part $\mathbf{B}$) as follows: $$x = A f(X) \textbackslash\{\}operatorname\{ch\} \textbackslash\{\}alpha T, \textbackslash\{\}qquad t = B  f(X)  \textbackslash\{\}operatorname\{sh\} \textbackslash\{\}alpha T, $$где $A$, $B$, $\alpha$ – постоянные, определённые в предыдущем пункте, а $f(X)$ – некоторая функция координаты $X$. Таким образом, пару чисел можно расcmатривать как координаты события в неинерциальной системе отсчета. Однако оказывается, что изменяя значения $T$, $X$ можно получить не все координаты $(t, x)$. Это связано с тем, что сигнал от некоторых событий никогда не достигает наблюдателя. Выразите функцию $f(X)$  через $g$ и $c$. Для простоты можете ограничиться случаем, когда событие произошло справа от наблюдателя. Если событие находится слева от наблюдателя, вид зависимости $(t, x)$ от $(T, X)$ такой же.  Изобразите в координатах $(x, t)$ область пространства-времени, которая покрывается координатами $(X,T)$, когда они меняются от $-\infty$ до $+\infty$.}{1.50}

\problembox{C3}{Recall that the interval between two infinitely close events is a quantity whose square is equal to $ds^2 = c^2 dt^2 - dx^2$. Show that the interval between two close events can be written as $$ ds^2= F(X, T) (c^2 dT^2 - dX^2). $$Найдите функцию $F(X, T)$. В ответ могут входить $g$, $c$.}{0.50}

\problembox{C4}{In the coordinate system constructed in this way, it is natural to consider bodies for which $X= \operatorname{const}$ to be motionless. Show that such bodies move with constant acceleration in the accompanying CO, express this acceleration $a(X)$ in terms of $X$, $g$, $c$. Thus, we can assume that the reference system is formed by observers who move according to the laws determined by the condition $X = \operatorname{const}$.}{0.70}

\problembox{C5}{Let an observer with coordinate $X_0$ at moment $T_0$ emit a light beam in the positive direction of the axis $x$. How will $X$ depend on $T$ for this ray at subsequent moments?}{0.70}

\problembox{C6}{Let an observer with coordinate $X_1$ at moment $T_1$ emit light with wavelength $\lambda_1$ (in the accompanying FR of the observer). Determine the wavelength $\lambda_2$ that will be measured by an observer with coordinate $X_2$ in his accompanying CO. Consider it $X_1, X_2 = \operatorname{const}$ for observers.}{1.30}

\vspace{1em}
\noindent\textit{Source:} \url{https://pho.rs/p/4563}

\end{document}